% Chapter 1: Introduction

\chapter{Introduction}
\label{ch:intro}

Disaster response operations place first responders in hazardous, unpredictable environments where both physical safety and decision-making under pressure are severely tested. Reconnaissance and situational awareness are essential for effective crisis management, yet they traditionally rely on personnel entering unstable structures and contested terrain. This project addresses these challenges by developing an autonomous Unmanned Ground Vehicle (UGV) designed to carry out reconnaissance while reducing human exposure and operator burden.

The primary aim is to improve the efficiency and safety of disaster response by shifting high-risk information-gathering tasks from humans to a capable robotic platform. The UGV is intended to perform mapping, hazard detection, and visual data collection in a coordinated manner, while supporting both centralized and decentralized communication. To move beyond raw data collection, the system incorporates AI-driven perception for automated situational assessment, allowing operators to receive interpreted information rather than unprocessed sensor streams.

A central technical objective is the implementation of LiDAR-based Simultaneous Localisation and Mapping (SLAM). This enables the UGV to estimate its own position reliably and to build accurate indoor maps of the environment, forming the foundation for autonomous navigation and for sharing spatial information with human operators and other agents. Data produced by the platform is conveyed through a multi-tier communication architecture that combines Wi-Fi, mesh networking, and long-range radio links, ensuring connectivity across varied and degraded conditions. To support extended or multi-robot missions, the UGV can autonomously deploy compact, waterproof mesh nodes that establish a resilient, decentralized network in environments where fixed infrastructure is absent or compromised.

Beyond mapping and connectivity, the system aims to reduce operator workload through an AI pipeline that analyses collected data to identify survivors, hazards, and navigable paths. The outputs are presented as human-readable reports and alerts, supporting faster and more informed decisions. The performance of the UGV is assessed in controlled simulations using defined Key Performance Indicators (KPIs), including localization drift, map accuracy, packet delivery ratio, network latency, detection accuracy, and the reduction in operator workload. These metrics provide a structured basis for evaluating the system against the stated goals of safer and more efficient disaster response.

This report presents the design, implementation, and evaluation of this autonomous indoor robot for disaster site exploration. The remainder of the document is organised as follows: the next chapter reviews related work; subsequent chapters describe the methodology, implementation, and results; and the report concludes with a summary and directions for future work.
